% Generated draft: requires mathematical and source review.
\documentclass[11pt]{article}
\usepackage[margin=25mm]{geometry}
\usepackage{fontspec}
\setmainfont{Latin Modern Roman}
\usepackage{amsmath,amssymb,mathtools,unicode-math}
\setmathfont{Latin Modern Math}
\usepackage{enumitem,needspace,xurl,hyperref}
\setlength{\parindent}{0pt}
\setlength{\parskip}{5pt}
\setlength{\emergencystretch}{3em}
\newcommand{\N}{\mathbb{N}}
\newcommand{\Z}{\mathbb{Z}}
\newcommand{\Q}{\mathbb{Q}}
\newcommand{\R}{\mathbb{R}}
\newcommand{\C}{\mathbb{C}}
\newcommand{\F}{\mathbb{F}}
\newcommand{\A}{\mathbb{A}}
\newcommand{\PP}{\mathbb{P}}
\newcommand{\E}{\mathbb{E}}
\newcommand{\eps}{\varepsilon}
\newcommand{\dd}{\,\mathrm{d}}
\begin{document}
\section*{UnsolvedMath ID: 30001656}
\textbf{Source problem code:} catalogue entry
\subsection{Definitions and mathematical statement}
A finite tournament is a directed graph \(T=(V(T),A(T))\) with no loops such that, for each pair of distinct vertices \(u,v\), exactly one of \((u,v)\) and \((v,u)\) belongs to \(A(T)\). For \(S\subseteq V(T)\), let \(T[S]\) denote the induced subtournament. A tournament is transitive if its vertices admit an ordering \(v_1,\ldots,v_m\) such that \((v_i,v_j)\in A(T)\) whenever \(i<j\). A tournament \(G\) is \(H\)-free if no induced subtournament of \(G\) is isomorphic to \(H\).

Define \(\chi(G)\) to be the least nonnegative integer \(k\) for which \(V(G)\) can be partitioned into \(k\) nonempty sets, each inducing a transitive subtournament; in particular, \(\chi(\varnothing)=0\).

Under the standard finite-tournament interpretation, the supplied Erdős--Hajnal conjecture for tournaments asks whether
\[
\forall H\;\exists\varepsilon_H\in\R\;\exists c_H>0\;\Bigl[\varepsilon_H<1\ \text{ and }\ \forall G\,\bigl((G\text{ is }H\text{-free and }|V(G)|\geq1)\Longrightarrow\chi(G)\leq c_H|V(G)|^{\varepsilon_H}\bigr)\Bigr],
\]
where \(H\) and \(G\) range over finite tournaments. Both constants may depend on \(H\), but neither may depend on \(G\).

The supplied statement does not explicitly specify finiteness or the treatment of the empty tournament. The displayed formulation makes the usual finite interpretation explicit and restricts the numerical inequality to nonempty \(G\), since the source permits \(\varepsilon_H\leq0\) and supplies no convention for \(0^{\varepsilon_H}\). No lower bound on \(\varepsilon_H\) has been added. The supplied background additionally describes a five-vertex regular tournament as an unresolved case, but gives no separate statement or explicit orientation for that example.
\subsection{Short English statement}
For each fixed forbidden tournament \(H\), the conjecture asks for constants \(c_H>0\) and \(\varepsilon_H<1\) such that every nonempty finite \(H\)-free tournament can be partitioned into at most \(c_H|V(G)|^{\varepsilon_H}\) transitive subtournaments. Finiteness and the empty-tournament boundary require the qualifications stated above.
\subsection{Sources}
\begin{itemize}
\item Local UnsolvedMath registry, supplied complete record with permanent ID \(30001656\), title ``Erdős--Hajnal Conjecture for Tournaments,'' exact source problem code \texttt{catalogue entry}; collection: Oberwolfach Reports Open Problems. The record gives proposed year 2011 and identifies no individual proposer.
\item Supplied report citation: \emph{Combinatorics} (2011), Oberwolfach Reports, DOI 10.4171/owr/2011/01, pp.\ 5--83. \url{https://doi.org/10.4171/owr/2011/01}. The supplied background describes this report as documenting workshop discussions and open problems.
\item Supplied literature source: ``Oberwolfach Report 01/2011,'' \url{https://doi.org/10.4171/owr/2011/01}. Literature checked at: \texttt{2026-08-21}. This is supplied dated curation, not an independently verified assessment. It labels the problem open and states that the Erdős--Hajnal reformulation remains open, including the five-vertex regular tournament named in the report. The background explicitly requests independent verification of time-sensitive status and cited claims.
\item The registry labels the statement \texttt{corrected\_verified} and reports that the official report was checked and the statement repaired to match its source passage. These are supplied provenance assertions, not verification performed for this draft.
\end{itemize}
\end{document}
